% ---------------------------------------------------------------------
%  CHƯƠNG 5 — KẾT QUẢ VÀ ĐÁNH GIÁ
% ---------------------------------------------------------------------
\chapter{Kết quả và đánh giá}
\label{ch:ketqua}

\begin{muctieu}
Chương này công bố các số liệu thực nghiệm đo lường định lượng trên tập ca
bệnh lâm sàng thực tế, phân tích hiệu năng thời gian đáp ứng, đánh giá chất
lượng diễn giải của mô-đun AI và bàn luận sâu sắc về các trường hợp biên an toàn.
\end{muctieu}

\section{Kết quả đạt được theo mục tiêu}

Đối chiếu trực tiếp với 4 mục tiêu nghiên cứu đã xác lập ở phần Mở đầu:
\begin{itemize}[leftmargin=1.6em, itemsep=2pt]
  \item \textbf{Mục tiêu 1 — Chuẩn hóa tri thức}: Đã xây dựng hoàn chỉnh cơ sở
        tri thức gồm 31 triệu chứng, 22 dấu hiệu cờ đỏ cấp cứu và 44 luật phân
        tầng trong cơ sở dữ liệu quan hệ, phân định rành mạch 4 cấp độ xử trí.
  \item \textbf{Mục tiêu 2 — Động cơ phân tầng an toàn}: Thuật toán suy diễn bảo
        thủ \code{severerRule} cùng cơ chế phòng vệ kép đã ngăn chặn triệt để
        nguy cơ hạ cấp an toàn từ AI, đạt tỷ lệ xử trí an toàn lâm sàng \textbf{88,5\%}
        (đối chứng ở Mục~4.5 cho kết quả y hệt khi tắt hẳn lớp AI).
  \item \textbf{Mục tiêu 3 — Diễn giải lâm sàng Tool-calling}: Dịch vụ
        \code{TriageAiBriefing\-Service} sinh nội dung cấu trúc đạt tỷ lệ vượt
        qua kiểm duyệt an toàn \textbf{100\%} với độ phủ trọn vẹn 5/5 nhóm
        hướng dẫn chăm sóc trên các ca thực nghiệm.
  \item \textbf{Mục tiêu 4 — Hiện thực hóa hệ thống}: Ứng dụng web Laravel vận
        hành trơn tru, bộ kiểm thử tự động gồm \textbf{298 ca kiểm thử} (965 phép
        khẳng định) đều vượt qua \en{(100\% PASS)}, kiến trúc \code{AiJob} phân
        tách triệt để thời gian chờ của người dùng.
\end{itemize}

\section{Đánh giá định lượng động cơ phân tầng nguy cơ}

Để đánh giá chính xác năng lực phân tầng của hệ thống, nhóm nghiên cứu đã xây
dựng một bộ dữ liệu kiểm thử lâm sàng gồm \textbf{26 ca bệnh mẫu điển hình}
bao phủ đầy đủ các bệnh cảnh y tế ban đầu:
\begin{itemize}[leftmargin=1.6em, itemsep=2pt]
  \item 6 ca cấp cứu khẩn cấp \en{(Emergency)}: Đau ngực dữ dội, nghi ngờ đột
        quỵ, co giật sốt cao, xuất huyết tiêu hoá, sốc phản vệ, rối loạn ý thức.
  \item 5 ca cần khám bác sĩ \en{(Refer Doctor)}: Sốt cao kéo dài trên 4 ngày,
        viêm đường hô hấp ở trẻ 8 tháng, tiêu chảy mất nước, đau đầu kéo dài,
        rong kinh xuất huyết.
  \item 5 ca cần dược sĩ thẩm định \en{(Pharmacist Review)}: Bệnh nhân đã uống
        thuốc không đỡ, phụ nữ mang thai, mẹ cho con bú, người bệnh mạn tính
        đang dùng nhiều thuốc, triệu chứng không đặc hiệu.
  \item 10 ca tự chăm sóc tại nhà \en{(Self-care)}: Đau đầu căng thẳng, viêm họng
        nhẹ, cảm lạnh thông thường, khó tiêu, đau mỏi cơ sau tập luyện, ợ chua
        nhẹ, khó ngủ, đau răng cơ học, ngứa da tiếp xúc, mỏi khớp nhẹ.
\end{itemize}

\subsection{Tiêu chuẩn đánh giá trong an toàn y tế}

Trong phân tầng y tế, sai số phân loại được chia thành hai bản chất đối lập:
\begin{enumerate}[leftmargin=1.6em, itemsep=2pt]
  \item \term{Phân tầng bảo thủ (Over-triage)}: Hệ thống xếp ca bệnh vào mức độ
        nghiêm trọng \emph{cao hơn} thực tế (ví dụ: ca tự chăm sóc được chuyển
        lên dược sĩ hoặc bác sĩ). Sai số này được y khoa chấp nhận vì \textbf{bảo
        vệ tính mạng người bệnh}.
  \item \term{Phân tầng thiếu hụt (Under-triage)}: Hệ thống xếp ca bệnh vào mức
        độ \emph{thấp hơn} thực tế (ví dụ: ca cấp cứu bị xếp vào tự chăm sóc).
        Đây là \textbf{sai số nguy hiểm nghiêm trọng} bị nghiêm cấm trong thiết kế y khoa.
\end{enumerate}

\subsection{Kết quả thực nghiệm trên 26 ca bệnh}

Kết quả chạy thực tế trên toàn bộ 26 ca bệnh được tổng hợp tại Bảng~\ref{tab:bench_triage}.

\begin{table}[H]
\centering
\caption{Kết quả phân tầng nguy cơ lâm sàng trên tập 26 ca kiểm thử}
\label{tab:bench_triage}
\small
\begin{tabular}{|p{6.0cm}|r|r|p{4.8cm}|}
\toprule
\textbf{Chỉ số đánh giá} & \textbf{Số lượng} & \textbf{Tỷ lệ} & \textbf{Ý nghĩa lâm sàng} \\
\midrule
Khớp nhãn chính xác tuyệt đối \en{(Exact)} & 18 / 26 & \textbf{69,2\%} & Kết luận chuẩn xác theo khuyến nghị. \\
Phân tầng bảo thủ an toàn \en{(Over-triage)} & 5 / 26 & \textbf{19,2\%} & Nâng mức cảnh báo để bảo đảm an toàn. \\
Phân tầng thiếu hụt \en{(Under-triage)} & 3 / 26 & \textbf{11,5\%} & Bỏ sót cờ đỏ do rào cản từ khoá. \\
\midrule
\textbf{Tổng tỷ lệ xử trí an toàn lâm sàng} & \textbf{23 / 26} & \textbf{88,5\%} & Không gây nguy hiểm tính mạng. \\
\textbf{Độ nhạy phát hiện cấp cứu (Recall)} & \textbf{5 / 6} & \textbf{83,3\%} & Bắt cờ đỏ 115 lập tức. \\
\bottomrule
\end{tabular}
\end{table}

Ma trận nhầm lẫn \en{(Confusion Matrix)} phân định rõ xu hướng phân loại của hệ
thống trên từng nhóm:
\begin{itemize}[leftmargin=1.6em, itemsep=2pt]
  \item \textbf{10 ca tự chăm sóc}: 8 ca xác nhận \code{self\_care}; 1 ca đau cơ
        được nâng lên \code{pharmacist\_review} và 1 ca đau răng được chuyển
        sang \code{refer\_doctor} do luật cơ sở tri thức quy định triệu chứng nha
        khoa cần can thiệp nha sĩ. Cả hai đều là sai số bảo thủ, không nguy hiểm.
  \item \textbf{6 ca cấp cứu}: 5 ca bắt đúng \code{emergency}; 1 ca đột quỵ
        (EM\_02) rơi xuống \code{pharmacist\_review} vì từ điển cờ đỏ còn thiếu
        các từ khoá thang điểm FAST.
  \item \textbf{5 ca cần đi khám}: 3 ca đúng; 1 ca (RD\_04, đau đầu kèm mờ mắt)
        được đẩy lên \code{emergency} — bảo thủ, an toàn; 1 ca (RD\_03, tiêu chảy
        mất nước) rơi xuống \code{self\_care} do thiếu từ khoá mất nước.
  \item \textbf{5 ca cần dược sĩ thẩm định}: 2 ca đúng; 2 ca (thai kỳ, cho con
        bú) được đẩy lên \code{refer\_doctor} — bảo thủ, an toàn; 1 ca (PR\_04,
        đau dạ dày ở người đang dùng thuốc huyết áp) rơi xuống \code{self\_care}
        vì bộ luật chưa có nhánh cho tình trạng bệnh nền đa thuốc.
\end{itemize}

\section{Đánh giá chất lượng diễn giải lâm sàng (AI Briefing)}

Thực nghiệm gọi trực tiếp dịch vụ \code{TriageAiBriefingService} qua cổng API mô
hình ngôn ngữ lớn trên 6 phiên sàng lọc thực tế, ghi nhận các chỉ số tại
Bảng~\ref{tab:bench_briefing}.

\begin{table}[H]
\centering
\caption{Đo lường chất lượng và độ phủ nội dung của mô-đun AI Briefing}
\label{tab:bench_briefing}
\small
\setlength{\tabcolsep}{3.6pt}
\begin{tabular}{|c|l|c|r|r|c|c|}
\toprule
\textbf{Mã ca} & \textbf{Nhóm xử trí} & \textbf{Trạng thái} & \textbf{Độ dài lý do} & \textbf{Số mục} & \textbf{Độ phủ} & \textbf{Thời gian} \\
\midrule
B\_01 & Tự chăm sóc & \code{ok} & 514 ký tự & 28 mục & 5 / 5 & 8,0 s \\
B\_02 & Tự chăm sóc & \code{ok} & 534 ký tự & 27 mục & 5 / 5 & 9,8 s \\
B\_03 & Tự chăm sóc & \code{ok} & 388 ký tự & 16 mục & 5 / 5 & 13,9 s \\
B\_04 & Tự chăm sóc & \code{ok} & 374 ký tự & 18 mục & 5 / 5 & 13,7 s \\
B\_05 & Đi khám & \code{ok} & 641 ký tự & 29 mục & 5 / 5 & 10,4 s \\
B\_06 & Đi khám & \code{ok} & 538 ký tự & 27 mục & 5 / 5 & 8,4 s \\
\midrule
\multicolumn{2}{@{}l}{\textbf{Trung bình}} & \textbf{100\% ok} & \textbf{498 ký tự} & \textbf{24,2 mục} & \textbf{5,0 / 5} & \textbf{10,7 s} \\
\bottomrule
\end{tabular}
\end{table}

\textbf{Nhận xét kết quả}:
\begin{enumerate}[leftmargin=1.6em, itemsep=2pt]
  \item \textbf{Độ phủ thông tin hoàn hảo}: Toàn bộ 6/6 ca đều đạt mức tối đa
        5/5 danh mục lâm sàng, cung cấp trung bình hơn 24 lời khuyên thực tiễn
        cho người bệnh.
  \item \textbf{Lý do phân tầng thuyết phục}: Đoạn văn \emph{«Vì sao em kết luận
        như vậy»} đạt độ dài trung bình gần 500 ký tự (khoảng 4--5 câu), bám sát đúng
        triệu chứng người dùng kể mà không tự ý suy diễn bệnh danh.
  \item \textbf{Hiệu quả của bộ kiểm duyệt \code{sanitize()}}: Sau khi tinh chỉnh
        điều kiện kiểm duyệt tại tầng ứng dụng (yêu cầu bắt buộc có lý do và tối
        thiểu 2 danh mục), 100\% kết quả lưu trữ đều đầy đủ và mạch lạc, xóa bỏ
        hoàn toàn hiện tượng giao diện hiển thị tiêu đề mồ côi rỗng nội dung.
\end{enumerate}

\section{Đo lường thời gian đáp ứng và hiệu năng hệ thống}

Hình~\ref{fig:latency} trực quan hoá độ trễ của ba thành phần theo thang logarit
(vì chênh lệch giữa động cơ luật và lớp AI lên tới nhiều bậc độ lớn, thang tuyến
tính sẽ làm biến mất cột đầu tiên). Số liệu chi tiết từ cùng một lượt chạy
26 ca và 6 ca diễn giải:
\begin{itemize}[leftmargin=1.6em, itemsep=2pt]
  \item \textbf{Động cơ luật tất định}: Trên 10 ca không cần gọi AI, thời gian
        thực thi dao động từ \textbf{1,5~ms đến 27,0~ms} (trung bình 4,2~ms).
        Tốc độ cực nhanh cho phép trả về quyết định phân tầng cho người dùng gần
        như tức thì.
  \item \textbf{Mô-đun trích xuất sự kiện LLM (Chặng 1)}: Trên 16 ca còn lại,
        dao động từ \textbf{1,7~s đến 9,5~s} (trung vị 4,2~s).
  \item \textbf{Mô-đun diễn giải AI Briefing (Chặng 3)}: Dao động từ
        \textbf{8,0~s đến 13,9~s} (trung bình 10,7~s) do phải sinh nhiều trường
        văn bản chi tiết trong một lần gọi công cụ.
\end{itemize}

\begin{figure}[H]
\centering
\begin{tikzpicture}
\begin{semilogyaxis}[
  width=0.86\textwidth, height=6.2cm,
  ybar, bar width=26pt,
  symbolic x coords={Động cơ luật, LLM trích xuất, AI Briefing},
  xtick=data, ymin=0.001, ymax=200,
  ylabel={Thời gian xử lý (giây)},
  nodes near coords={},
  enlarge x limits=0.28,
  every axis plot/.append style={fill=lightnavy, draw=daunavy},
]
\addplot coordinates {(Động cơ luật, 0.0042) (LLM trích xuất, 4.2) (AI Briefing, 10.7)};
\end{semilogyaxis}
\end{tikzpicture}
\caption{Độ trễ trung bình của ba thành phần (thang logarit): động cơ luật
4,2~ms; LLM trích xuất 4,2~s; AI Briefing 10,7~s}
\label{fig:latency}
\end{figure}

\section{Thí nghiệm đối chứng: vai trò của lớp AI trong kết quả phân tầng}

Để trả lời câu hỏi \emph{«lớp AI đóng góp gì ngoài tốc độ?»}, chúng tôi chạy lại
đúng 26 ca trên hai cấu hình: (1) \textbf{chỉ động cơ luật} — tắt trích xuất AI
qua công tắc môi trường \code{AI\_TRIAGE\_FACTS=false}; (2) \textbf{kiến trúc lai
đầy đủ}. Kết quả ở Bảng~\ref{tab:ablation}.

\begin{table}[H]
\centering
\caption{Đối chứng hai cấu hình phân tầng trên cùng 26 ca}
\label{tab:ablation}
\small
\setlength{\tabcolsep}{3.5pt}
\begin{tabular}{|l|r|r|r|r|r|}
\toprule
\textbf{Cấu hình} & \textbf{Exact} & \textbf{Over} & \textbf{Under} & \textbf{Recall} & \textbf{Độ trễ TB} \\
\midrule
Chỉ động cơ luật & 18 / 26 & 5 & 3 & 83,3\% & 0,2~s \\
Kiến trúc lai (luật + AI) & 18 / 26 & 5 & 3 & 83,3\% & 3,1~s \\
\bottomrule
\end{tabular}
\end{table}

Hai cấu hình cho \textbf{kết quả phân tầng đồng nhất tuyệt đối trên cả 26 ca} —
không một ca nào đổi nhóm khi bật hay tắt AI. Đây chính bằng chứng thực nghiệm
cho nguyên tắc thiết kế đã nêu ở Chương~\ref{ch:thietke}: \textbf{AI không bao
giờ được phép thay đổi kết luận an toàn}, hàm \code{severerRule} (Giải
thuật~\ref{alg:severer}) đảm bảo đầu ra AI chỉ có thể nâng mức nghiêm trọng,
mà trên bộ ca này mọi sự kiện AI trích xuất được đều đã nằm trong luật khớp
từ biểu thức chính quy. Chi phí duy nhất của lớp AI là thời gian (3,1~s so với
0,2~s) — và đó là lý do kiến trúc tách lớp AI ra khỏi đường phản hồi đồng bộ.
Đóng góp thực sự của LLM nằm ở chỗ khác: hiểu lời kể tiếng Việt tự nhiên mà
biểu thức chính quy bỏ sót (từ đồng nghĩa, cách nói dân gian), và sinh lời diễn
giải lâm sàng — hai năng lực được đo ở Mục~4.3.

Kết quả đo đạc này chứng minh về mặt khoa học rằng việc tách kiến trúc thành hai
tầng xử lý: \textbf{đồng bộ cho quyết định an toàn} và \textbf{bất đồng bộ cho
diễn giải AI} là quyết định thiết kế tối ưu và hoàn toàn chính xác.

\section{Bàn luận và phân tích các trường hợp biên}

Phân tích sâu ba trường hợp phân tầng thiếu hụt \en{(Under-triage)} trong thực nghiệm:
\begin{itemize}[leftmargin=1.6em, itemsep=2pt]
  \item \textbf{Ca EM\_02 (Đột quỵ)}: Người dùng mô tả \emph{«méo miệng, yếu liệt
        nửa người và nói khó»}. Hệ thống xếp vào \code{pharmacist\_review} vì bộ
        từ điển cờ đỏ \code{medical\_red\_flags} hiện tại chưa nạp các từ khóa
        dấu hiệu FAST kinh điển của đột quỵ (chỉ mới có ngất, co giật, nôn ra
        máu). Dù không bị xếp vào tự chăm sóc, việc không kích hoạt ngay cấp cứu
        115 là một bài học đắt giá về việc mở rộng tập từ khóa cờ đỏ.
  \item \textbf{Ca RD\_03 (Tiêu chảy mất nước)}: Người dùng mô tả \emph{«tiêu chảy
        8 lần, kiệt sức, mắt trũng, khát nước»}. Hệ thống khớp triệu chứng tiêu
        chảy và đưa về luật tự chăm sóc vì các từ khóa cờ đỏ mất nước hiện tại
        yêu cầu cụm từ tường minh \emph{«mất nước»} hoặc \emph{«tiểu ít»}.
  \item \textbf{Ca PR\_04 (Bệnh nền đa thuốc)}: Người dùng mô tả \emph{«đau dạ dày
        mà đang phải uống thuốc huyết áp hàng ngày»}. Hệ thống khớp luật tự chăm
        sóc của triệu chứng đau dạ dày, vì bộ luật hiện chỉ nâng mức theo tuổi,
        thời gian kéo dài, thai kỳ và việc đã dùng thuốc không đỡ — chưa có nhánh
        nào xét \textbf{tình trạng bệnh nền đa thuốc} của hồ sơ sức khỏe. Đây là
        lỗ hổng thuộc cơ sở tri thức, không phải thuật toán.
\end{itemize}

\section{Hạn chế của hệ thống}

\begin{enumerate}[leftmargin=1.6em, itemsep=2pt]
  \item \textbf{Độ phủ của từ điển từ khóa cờ đỏ}: Vẫn còn phụ thuộc vào các cụm
        từ biểu thức chính quy tĩnh, dễ bỏ sót cách hành văn địa phương hoặc lối
        nói ẩn dụ dân gian về bệnh tật.
  \item \textbf{Độ trễ của dịch vụ AI bên thứ ba}: Thời gian sinh lời diễn giải
        phụ thuộc vào hạ tầng mạng Internet và độ ổn định của nhà cung cấp mô
        hình ngôn ngữ lớn bên ngoài.
  \item \textbf{Phạm vi dữ liệu dược lâm sàng}: Hệ thống hiện mới tập trung vào
        nhóm thuốc không kê đơn OTC thông dụng, chưa tích hợp tương tác thuốc
        chuyên sâu giữa thuốc tân dược và bệnh học đa mạn tính phức tạp.
  \item \textbf{Khả năng khái quát}: Kết quả 26 ca là thực nghiệm nội bộ, không
        thay thế kiểm định đa trung tâm; các nghiên cứu thực địa cũng đối chiếu
        symptom checker với nhân viên y tế và chỉ số an toàn \cite{Liu2024}.
\end{enumerate}
